\section{结论}

MRI由形态显示走向决策支持，是信号编码对象、定量模型与图像分析方法共同演进的结果。DWI/ADC把扩散过程带进临床，DTI和DKI接着描述方向性与非高斯微结构；DCE-MRI靠动态采集和药代动力学模型，让强化图像承载灌注与通透性信息；BOLD-fMRI借助快速成像和神经血管耦合，把观察对象由局部脑活动扩展到静息态脑网络。这些技术都没有取代常规MRI，它们在其解剖基础上加了若干维度的证据。

多参数融合、影像组学和人工智能把病理表型与风险预测又往前推了一步，但案例同样表明，联合模型并非总是显著优于简单模型，高AUC也不会自动变成普适的临床效用。写到最后，我的认识是这样：MRI技术发展的目标不是“给出最多信息”，而是“以可接受的代价给出最可靠、与决策相关的那部分信息”。值得关注的方向，是快速采集、定量模型、规范流程与临床验证一起往前走。对生物医学工程专业的学生来说，把这条从物理信号到患者获益的证据链读懂，并且动手检验它，比多背几个技术名词有用。
